\documentclass[10pt,a4paper]{amsart}
\usepackage[margin=19mm]{geometry}
\usepackage{amsmath,amssymb,mathtools}
\usepackage[scr=rsfs,cal=euler]{mathalfa}
\usepackage{xfrac}
\usepackage[unicode,hidelinks,bookmarks=false]{hyperref}
\newcommand{\ie}{i.e.\ }
\newcommand{\cf}{cf.\ }
\newcommand{\resp}{respectively\ }
\newcommand{\Subsection}{\subsection}
\providecommand{\nobreakdash}{\nobreak\hbox{-}}
\setlength{\emergencystretch}{2em}
\multlinegap0pt
\usepackage{bm}
\usepackage{bbm}
\usepackage{dsfont}
\usepackage{xspace}
\usepackage{cancel}
\usepackage{mathtools}
\usepackage{amsthm}
\makeatletter
\@ifundefined{theorem}{\newtheorem{theorem}{Theorem}}{}
\@ifundefined{proposition}{\newtheorem{proposition}[theorem]{Proposition}}{}
\makeatother
\def\PP{{\mathbb P}}
\def\Tscr{{\mathscr T}}
\def\Uscr{{\mathscr U}}
\def\ZZ{{\mathbb Z}} 
\def\g{\gamma}
\def\la{\langle}
\def\ra{\rangle}
\newcommand\MW{\operatorname{MW}}
\newcommand\Trans{\operatorname{Trans}}
\newcommand{\eps}{\varepsilon}
\newcommand{\p}{\partial}
\title[The smooth Mordell-Weil group and mapping class groups of el]{The smooth Mordell-Weil group and mapping class groups of elliptic surfaces}
\author{Benson Farb, Eduard Looijenga}
\date{Source: arXiv:2403.15960v3 (CC BY 4.0)}
\begin{document}
\maketitle

\noindent\emph{Source and licence.} The smooth Mordell-Weil group and mapping class groups of elliptic surfaces. Version arXiv:2403.15960v3 (CC BY 4.0), distributed under the Creative Commons Attribution 4.0 International licence, as recorded in the arXiv licence metadata for this exact version and shown on the abstract page https://arxiv.org/abs/2403.15960v3. Full rights notice: licenses/arxiv-2403.15960-CC-BY-4.0.txt.\par\smallskip

\noindent\textit{Editorial scope.} Counted unit: the Proposition whose source label is \texttt{thm:MWcohomologyI} in the article (source0.tex lines 932--941, source label \texttt{thm:MWcohomologyI}), delivered together with the article's own complete original proof of it (source0.tex lines 978--1043, one \texttt{proof} environment of 66 lines). No mathematics is written, repaired or re-derived: statement and proof are the article's own text line for line. The proof is detached from the statement in the source: the article states the result at lines 932--941 and prints its proof at lines 978--1043, and the proof environment's own header names the statement, so the association is the article's own. Required context retained uncounted: prop:MWcohomology (lines 868--879); prop:MWcohomologyII (lines 943--960). Every definition, display and label of those lines is retained unchanged. Omitted from this package and not redistributed: 1--867, 880--931, 942--942, 961--977, 1044--1839. Nothing inside a retained excerpt is omitted or altered: every retained body is delivered exactly as the article's lines are written, so no omission receipt of any kind accompanies this package. Citations: the retained excerpts carry no citation command at all, so no bibliography entry of the article is transcribed and no reference is redistributed. Reference adaptation: none was needed and none was made; every reference inside the delivered unit resolves inside it, so no unresolved reference or citation is produced. Printed slips kept literal: no printed word, symbol, hypothesis or number of the counted statement or of its proof was altered. Layout and class: the article's own class is \texttt{amsart} with packages including amsmath, amsthm, amsbsy, amsfonts, amssymb, txfonts, stmaryrd, mathrsfs, graphicx, amscd, tikz-cd, tikz, cancel, enumitem; the wrapper is the lane's pinned \texttt{amsart} editorial wrapper at 10pt on A4, which restates the theorem-like environments the retained text needs and adds the article's own macro definitions that the retained text needs (\texttt{\textbackslash MW}, \texttt{\textbackslash PP}, \texttt{\textbackslash Trans}, \texttt{\textbackslash Tscr}, \texttt{\textbackslash Uscr}, \texttt{\textbackslash ZZ}, \texttt{\textbackslash eps}, \texttt{\textbackslash g}, \texttt{\textbackslash la}, \texttt{\textbackslash p}, \texttt{\textbackslash ra}), with extra wrapper packages (bm, bbm, dsfont, xspace, cancel, mathtools). The delivered numbering is the unit's own. Assets: no retained span contains a figure environment, a graphics-inclusion command, a PDF-page inclusion, a table or a TikZ picture; no image file is copied into this package, the package declares no asset dependency and no image file is redistributed with it. Licence: version arXiv:2403.15960v3 of this article is distributed under the Creative Commons Attribution 4.0 International licence, as recorded in the arXiv licence metadata for this exact version and shown on the abstract page https://arxiv.org/abs/2403.15960v3. A full rights notice is delivered at licenses/arxiv-2403.15960-CC-BY-4.0.txt. Characters: every retained excerpt is pure ASCII, so the delivered engine, XeLaTeX with its default Latin Modern text font, reports no missing character.
\begin{proposition}[{\bf Cohomological characterization of MW groups}]
\label{prop:MWcohomology}
 Assume that $\pi$  has only integral fibers  (all singular  fibers are either nodal or cuspidal). Then there are natural isomorphisms 
 \[
 \MW(\pi)\cong H^1(B, R^1\pi_*\ZZ) \ \ \text{and}\ \  \MW(\pi, \p\pi) \cong H^1(B, \p B; R^1\pi_*\ZZ).
 \]
 Moreover,  for $q>0$, there are natural isomorphisms 
\[
H^q(B, \Tscr(\pi))\cong  H^{q+1}(B,R^1\pi_*\ZZ), \quad 
H^q(B,\p B;\Tscr(\pi))\cong  H^{q+1}(B,D_\pi;R^1\pi_*\ZZ).
\]
\end{proposition}

\begin{proposition}[{\bf Existence of sections and action of $\MW(\pi)$ II}]
\label{prop:MWcohomologyII}
Let  $\pi :M\to \PP^1$ be a genus one fibration of arithmetic  genus $d\ge 1$ with integral fibers (so $M\cong M_d$).
Then $\pi$ admits a section and for every  $\tau\in \MW(\pi)$ there exists a unique $c_\tau\in e^\perp/\ZZ e$
(where  $e\in H_2(M)$ is the fiber class)
 such that $\tau$ acts on $H_2(M)$ as the Eichler transformation $E(e \wedge c_\tau)$ given by 
\[
E(e\wedge c_\tau)(x)=x+ (x \cdot e )\hat c_\tau - (x\cdot \hat c_\tau)e -\tfrac{1}{2}(c_\tau\cdot c_\tau)(x\cdot e)e \]
for all $x\in H_2(M)$, 
where $\hat c_\tau\in e ^\perp$ is a lift of $c_\tau$. 
The resulting map 
\[
\tau\in\MW(\pi)\to  c_\tau\in e^\perp/\ZZ e
\]
 is an isomorphism of abelian groups. In particular, $\MW(\pi)$ has  naturally the structure of a lattice (which is even unimodular and of 
 signature $(2d-2, 10d-2)$). 

\end{proposition}

\begin{proposition}[{\bf Existence of sections and action of $\MW(\pi)$ I}]
\label{thm:MWcohomologyI}
Let  $\pi :X\to B$ be a genus one fibration over  a compact, connected, oriented surface $B$ with nonempty boundary 
$\p B$.  Assume that each fiber of $\pi$ is integral.  Denote by $e\in H_2(X)$ the homology class of a general fiber and  by  $H^2(X)^e$ its annihilator 
in $H^2(X)$ (i.e., the cohomology classes that vanish on $e$). 

   Then $\pi$ admits a section. For every  $\tau\in \MW(\pi)$ there exists a unique $\g_\tau\in H^2(X)$ such that for all $\xi\in H^2(X)$: 
\[\tau(\xi)= \xi+\la \xi|e\ra \g_\tau\]
and the resulting map $\tau\in \MW(\pi)\mapsto  \g_\tau\in H^2(X)^e$ is an isomorphism of abelian groups.
\end{proposition}

\begin{proof}[Proof of  Propositions  \ref{thm:MWcohomologyI} and \ref{prop:MWcohomologyII}]
We prove the propositions  in a number of steps.
Our assumption that the fibers of $\pi$ are integral implies that 
$R^2\pi_*\ZZ$ is  a trivial local system of rank one on $B$ so that a natural isomorphism  $H^0(B, R^2\pi_*\ZZ)\cong \ZZ 
$ is given by integration over (=natural pairing with) the fiber class, i.e., the map $\xi\in H^2(X)\mapsto \la \xi|e\ra \in \ZZ$.
\\

\emph{Step 1: The Mordell-Weil group $\MW(\pi)$ acts trivially on the grading of $H^2(X)$ with respect to  the Leray filtration defined by $\pi$. }

Consider the Leray spectral sequence
\[
E_2^{p,q}= H^p(B, R^q\pi_*\ZZ)\Rightarrow H^{p+q}(X).
\]  
Any diffeomorphism of $X$ that is a fiberwise-translation acts trivially on each term $E_2^{p,q}$. Hence the same is true for  its action  on $E_\infty^{p,q}$, which is by definition the grading of $H^2(X)$ with respect to  the Leray filtration.  
\\

\emph{Step 2: If $H^{2}(B,R^1\pi_*\ZZ)$ vanishes then $\pi$ admits a section.}

Since $\pi$ has no multiple fibers, we can find  an open cover $\Uscr$ of  $B$ such that  for each $U\in \Uscr$, $\pi_{U}$ admits a section $s_U$. Then $s_{U'}-s_U$ defines a section of $\Trans(\pi)$ over $U\cap U'$. Together they make up a \v{C}ech $1$-cycle relative the covering $\Uscr$ and hence define an element of $\check{H}^1(\Uscr,\Trans(\pi))$. If this a coboundary, then
there exist $\{t_U\in H^0(U, \Trans(\pi))\}_{U\in \Uscr}$ such that $s_{U'}-s_U=t_{U'}-t_U$, which means that the collection $\{s_U+t_U\}_{U\in\Uscr}$ defines a global section. Since we are allowed to pass to a refinement of $\Uscr$, it suffices to show that $\check{H}^1(B,\Trans(\pi))$ is trivial. For paracompact spaces \v{C}ech cohomology is ordinary cohomology, and so this  is  by Proposition \ref{prop:MWcohomology} equal to $H^{2}(B,R^1\pi_*\ZZ)$. 
Hence Step 2 follows.
\\

\emph{Step 3: Proof  of  Proposition \ref{thm:MWcohomologyI}.}

Then $B$ contains an embedded  graph $G\subset B$ as a deformation retract which  contains all the  critical values. This implies that the restriction map  $E_2^{p,q}= H^p(B, R^q\pi_*\ZZ)\to H^p(G, R^q\pi_*\ZZ)$ is an isomorphism  so that $E_2^{p,q}$ is trivial unless 
$p\in \{0,1\}$ and $q\in \{0,1,2\}$.  In particular, $\pi$  admits a section by Step 2. It also follows that  the Leray sequence degenerates on this page, so that  there is an exact sequence
\[
0\to H^1(B, R^1\pi_*\ZZ)\to H^2(X)\to H^0(B, R^2\pi_*\ZZ)\to 0.
\]
The  map $H^2(X)\to H^0(B, R^2\pi_*\ZZ)\cong \ZZ$ is given by integration over the fiber class $e$. It is 
surjective, because it takes the value $1$ on a section (we here regard the class of a section as an element of 
$H_2(X, \p X)\cong H^2(X)$). Proposition \ref{prop:MWcohomology}  identifies  $\MW(\pi)$ with its kernel of 
this map.  Any  $\tau\in \MW(\pi)$ induces a transformation of $H^2(X)$ which preserves the above short 
exact sequence  and acts trivially on both $H^1(B, R^1\pi_*\ZZ)$ and $H^0(B, R^2\pi_*\ZZ)\cong\ZZ$. 
Hence it is of the form $\xi\in H^2(X)\mapsto \xi +\la \xi\, |\, e\ra \g_\tau$ for some 
$\g_\tau\in H^1(B, R^1\pi_*\ZZ)$. So if  $\xi$ is the class  of a section, then its image under $\tau$ is $\xi+\g_\tau$. 
The definition shows that $\tau\mapsto \g_\tau$ is in fact the isomorphism 
$\MW(\pi)\cong H^1(B, R^1\pi_*\ZZ)$ found above.
\\

\emph{Step 4: Proof  of  Proposition \ref{prop:MWcohomologyII}}

We first show that the differential $\ZZ\cong H^0(\PP^1, R^2\pi_*\ZZ)\to H^2(\PP^1, R^1\pi_*\ZZ)$ in the Leray 
spectral sequence is the zero map. Indeed, we already observed that integration over $e$ defines 
an isomorphism $H^0(\PP^1, R^2\pi_*\ZZ)\cong \ZZ$.
Since the Leray spectral sequence degenerates on its third page, the kernel of this differential is  a 
quotient of $H^2(M)$ and it is via this quotient that the integration map 
$\xi\in H^2(M)\mapsto \la\xi | e\ra$ is defined. Since $e$ is primitive, the latter is onto and hence this kernel 
is all of  $H^0(\PP^1, R^2\pi_*\ZZ)$. So the differential $H^0(\PP^1, R^2\pi_*\ZZ)\to H^2(\PP^1, R^1\pi_*\ZZ)$ must be \
the zero map. 

This implies that $H^2(\PP^1, R^1\pi_*\ZZ)$ embeds in $H^3(M)$. But $H^3(M)\cong H_1(M)=0$ and hence $H^2(\PP^1, R^1\pi_*\ZZ)=0$.
So $\pi$ admits a section by Step 2. This in turn implies that the natural map $H^2(\PP^1)\to H^2(M)$ is injective.
Hence the Leray filtration  of  $H^2(M)$ is given by 
\[
H^2(M)=L_2H^2\supset L_1H^2\supset L_0H^2=H^2(\PP^1)\supset L_{-1}H^2=0
\]
with $L_1H^2/L_0H^2\cong H^1(\PP^1, R^1\pi_*\ZZ)$ and $L_2H^2/L_1H^2\cong  H^0(\PP^1, R^2\pi_*\ZZ)$. 
The preceding shows that $L_1H^2$ is the annihilator of $e$ and that $L_0H^2=H^2(\PP^1)$ is spanned by the Poincar\'e dual $\eps $ of $e$. If $\sigma\in H_2(M)$ is the class of  a section, then its Poincar\'e dual and $\eps $ span a unimodular lattice of rank $2$  of signature  $(1,1)$ (even or odd according the parity of $\sigma\cdot\sigma=-d$).  
Since the intersection pairing on $H_2(M)$ is unimodular, its orthogonal complement is also unimodular. This 
orthogonal complement maps isomorphically onto $\eps ^\perp/\ZZ \eps $ and hence the latter  is unimodular. 

Any orthogonal transformation of the unimodular lattice  $H_2(M)$ whose induced transformation on $H^2(M)$ acts trivially on the successive quotients of the above filtration (defined by the isotropic vector $\eps$) is an Eichler transformation as stated in the theorem  with $c$ canonically defined in $\Lambda_d(e)$. In particular, we find for every $\tau\in \MW(\pi)$ an element $c_\tau\in \Lambda_d(e)$. For $\sigma$  as above,  $\tau (\sigma)\equiv \sigma+c_\tau\pmod{\ZZ \eps }$.
Again the  definition shows that $\tau\mapsto c_\tau$ is in fact the isomorphism $\MW(\pi)\cong H^1(\PP^1, R^1\pi_*\ZZ)\cong e^\perp/\ZZ e$. 
\end{proof}

\end{document}
